% Options for packages loaded elsewhere
\PassOptionsToPackage{unicode}{hyperref}
\PassOptionsToPackage{hyphens}{url}
\documentclass[
]{article}
\usepackage{xcolor}
\usepackage[margin=1in]{geometry}
\usepackage{amsmath,amssymb}
\setcounter{secnumdepth}{-\maxdimen} % remove section numbering
\usepackage{iftex}
\ifPDFTeX
  \usepackage[T1]{fontenc}
  \usepackage[utf8]{inputenc}
  \usepackage{textcomp} % provide euro and other symbols
\else % if luatex or xetex
  \usepackage{unicode-math} % this also loads fontspec
  \defaultfontfeatures{Scale=MatchLowercase}
  \defaultfontfeatures[\rmfamily]{Ligatures=TeX,Scale=1}
\fi
\usepackage{lmodern}
\ifPDFTeX\else
  % xetex/luatex font selection
\fi
% Use upquote if available, for straight quotes in verbatim environments
\IfFileExists{upquote.sty}{\usepackage{upquote}}{}
\IfFileExists{microtype.sty}{% use microtype if available
  \usepackage[]{microtype}
  \UseMicrotypeSet[protrusion]{basicmath} % disable protrusion for tt fonts
}{}
\makeatletter
\@ifundefined{KOMAClassName}{% if non-KOMA class
  \IfFileExists{parskip.sty}{%
    \usepackage{parskip}
  }{% else
    \setlength{\parindent}{0pt}
    \setlength{\parskip}{6pt plus 2pt minus 1pt}}
}{% if KOMA class
  \KOMAoptions{parskip=half}}
\makeatother
\usepackage{color}
\usepackage{fancyvrb}
\newcommand{\VerbBar}{|}
\newcommand{\VERB}{\Verb[commandchars=\\\{\}]}
\DefineVerbatimEnvironment{Highlighting}{Verbatim}{commandchars=\\\{\}}
% Add ',fontsize=\small' for more characters per line
\usepackage{framed}
\definecolor{shadecolor}{RGB}{248,248,248}
\newenvironment{Shaded}{\begin{snugshade}}{\end{snugshade}}
\newcommand{\AlertTok}[1]{\textcolor[rgb]{0.94,0.16,0.16}{#1}}
\newcommand{\AnnotationTok}[1]{\textcolor[rgb]{0.56,0.35,0.01}{\textbf{\textit{#1}}}}
\newcommand{\AttributeTok}[1]{\textcolor[rgb]{0.13,0.29,0.53}{#1}}
\newcommand{\BaseNTok}[1]{\textcolor[rgb]{0.00,0.00,0.81}{#1}}
\newcommand{\BuiltInTok}[1]{#1}
\newcommand{\CharTok}[1]{\textcolor[rgb]{0.31,0.60,0.02}{#1}}
\newcommand{\CommentTok}[1]{\textcolor[rgb]{0.56,0.35,0.01}{\textit{#1}}}
\newcommand{\CommentVarTok}[1]{\textcolor[rgb]{0.56,0.35,0.01}{\textbf{\textit{#1}}}}
\newcommand{\ConstantTok}[1]{\textcolor[rgb]{0.56,0.35,0.01}{#1}}
\newcommand{\ControlFlowTok}[1]{\textcolor[rgb]{0.13,0.29,0.53}{\textbf{#1}}}
\newcommand{\DataTypeTok}[1]{\textcolor[rgb]{0.13,0.29,0.53}{#1}}
\newcommand{\DecValTok}[1]{\textcolor[rgb]{0.00,0.00,0.81}{#1}}
\newcommand{\DocumentationTok}[1]{\textcolor[rgb]{0.56,0.35,0.01}{\textbf{\textit{#1}}}}
\newcommand{\ErrorTok}[1]{\textcolor[rgb]{0.64,0.00,0.00}{\textbf{#1}}}
\newcommand{\ExtensionTok}[1]{#1}
\newcommand{\FloatTok}[1]{\textcolor[rgb]{0.00,0.00,0.81}{#1}}
\newcommand{\FunctionTok}[1]{\textcolor[rgb]{0.13,0.29,0.53}{\textbf{#1}}}
\newcommand{\ImportTok}[1]{#1}
\newcommand{\InformationTok}[1]{\textcolor[rgb]{0.56,0.35,0.01}{\textbf{\textit{#1}}}}
\newcommand{\KeywordTok}[1]{\textcolor[rgb]{0.13,0.29,0.53}{\textbf{#1}}}
\newcommand{\NormalTok}[1]{#1}
\newcommand{\OperatorTok}[1]{\textcolor[rgb]{0.81,0.36,0.00}{\textbf{#1}}}
\newcommand{\OtherTok}[1]{\textcolor[rgb]{0.56,0.35,0.01}{#1}}
\newcommand{\PreprocessorTok}[1]{\textcolor[rgb]{0.56,0.35,0.01}{\textit{#1}}}
\newcommand{\RegionMarkerTok}[1]{#1}
\newcommand{\SpecialCharTok}[1]{\textcolor[rgb]{0.81,0.36,0.00}{\textbf{#1}}}
\newcommand{\SpecialStringTok}[1]{\textcolor[rgb]{0.31,0.60,0.02}{#1}}
\newcommand{\StringTok}[1]{\textcolor[rgb]{0.31,0.60,0.02}{#1}}
\newcommand{\VariableTok}[1]{\textcolor[rgb]{0.00,0.00,0.00}{#1}}
\newcommand{\VerbatimStringTok}[1]{\textcolor[rgb]{0.31,0.60,0.02}{#1}}
\newcommand{\WarningTok}[1]{\textcolor[rgb]{0.56,0.35,0.01}{\textbf{\textit{#1}}}}
\usepackage{longtable,booktabs,array}
\newcounter{none} % for unnumbered tables
\usepackage{calc} % for calculating minipage widths
% Correct order of tables after \paragraph or \subparagraph
\usepackage{etoolbox}
\makeatletter
\patchcmd\longtable{\par}{\if@noskipsec\mbox{}\fi\par}{}{}
\makeatother
% Allow footnotes in longtable head/foot
\IfFileExists{footnotehyper.sty}{\usepackage{footnotehyper}}{\usepackage{footnote}}
\makesavenoteenv{longtable}
\usepackage{graphicx}
\makeatletter
\newsavebox\pandoc@box
\newcommand*\pandocbounded[1]{% scales image to fit in text height/width
  \sbox\pandoc@box{#1}%
  \Gscale@div\@tempa{\textheight}{\dimexpr\ht\pandoc@box+\dp\pandoc@box\relax}%
  \Gscale@div\@tempb{\linewidth}{\wd\pandoc@box}%
  \ifdim\@tempb\p@<\@tempa\p@\let\@tempa\@tempb\fi% select the smaller of both
  \ifdim\@tempa\p@<\p@\scalebox{\@tempa}{\usebox\pandoc@box}%
  \else\usebox{\pandoc@box}%
  \fi%
}
% Set default figure placement to htbp
\def\fps@figure{htbp}
\makeatother
\setlength{\emergencystretch}{3em} % prevent overfull lines
\providecommand{\tightlist}{%
  \setlength{\itemsep}{0pt}\setlength{\parskip}{0pt}}
\usepackage{bookmark}
\IfFileExists{xurl.sty}{\usepackage{xurl}}{} % add URL line breaks if available
\urlstyle{same}
\hypersetup{
  hidelinks,
  pdfcreator={LaTeX via pandoc}}

\author{}
\date{\vspace{-2.5em}}

\begin{document}

\section{smMIP Analysis Pipeline
Overview}\label{smmip-analysis-pipeline-overview}

\subsection{Quick Start}\label{quick-start}

\begin{Shaded}
\begin{Highlighting}[]
\FunctionTok{bash}\NormalTok{ /Volumes/Seq\_SSD/smMIP/Universal/Code/Master\_pipeline.sh }\DataTypeTok{\textbackslash{}}
  \AttributeTok{{-}{-}directory}\NormalTok{ /Volumes/Seq\_SSD/smMIP/KG001\_01.22.25 }\DataTypeTok{\textbackslash{}}
  \AttributeTok{{-}{-}fastq\_dir}\NormalTok{ /Volumes/Seq\_SSD/smMIP/KG001\_01.22.25/RAW\_FASTQ}
\end{Highlighting}
\end{Shaded}

\subsection{Pipeline Stages}\label{pipeline-stages}

{\def\LTcaptype{none} % do not increment counter
\begin{longtable}[]{@{}
  >{\raggedright\arraybackslash}p{(\linewidth - 4\tabcolsep) * \real{0.2500}}
  >{\raggedright\arraybackslash}p{(\linewidth - 4\tabcolsep) * \real{0.2857}}
  >{\raggedright\arraybackslash}p{(\linewidth - 4\tabcolsep) * \real{0.4643}}@{}}
\toprule\noalign{}
\begin{minipage}[b]{\linewidth}\raggedright
Stage
\end{minipage} & \begin{minipage}[b]{\linewidth}\raggedright
Script
\end{minipage} & \begin{minipage}[b]{\linewidth}\raggedright
Description
\end{minipage} \\
\midrule\noalign{}
\endhead
\bottomrule\noalign{}
\endlastfoot
1 & BWA.sh & Align paired-end reads to hg19 reference \\
2 & filter\_bam.sh & Filter for quality (MAPQ\textgreater=50, properly
paired) \\
3 & Read\_Processing.sh & Map reads to smMIP probes, extract UMIs \\
4 & Level\_Bas\_Calls.sh & Generate variant pileups at smMIP level \\
5 & Annotate\_SNVs.R & Annotate panel with gene/protein/COSMIC info \\
6 & calling\_mutations.R & Call mutations using statistical error
modeling \\
\end{longtable}
}

\subsection{Command-Line Parameters}\label{command-line-parameters}

\subsubsection{Required Parameters}\label{required-parameters}

{\def\LTcaptype{none} % do not increment counter
\begin{longtable}[]{@{}ll@{}}
\toprule\noalign{}
Parameter & Description \\
\midrule\noalign{}
\endhead
\bottomrule\noalign{}
\endlastfoot
\texttt{-\/-directory} & Base experiment directory (outputs go here) \\
\texttt{-\/-fastq\_dir} & Directory containing raw FASTQ files \\
\end{longtable}
}

\subsubsection{Optional Parameters}\label{optional-parameters}

{\def\LTcaptype{none} % do not increment counter
\begin{longtable}[]{@{}
  >{\raggedright\arraybackslash}p{(\linewidth - 4\tabcolsep) * \real{0.3333}}
  >{\raggedright\arraybackslash}p{(\linewidth - 4\tabcolsep) * \real{0.2727}}
  >{\raggedright\arraybackslash}p{(\linewidth - 4\tabcolsep) * \real{0.3939}}@{}}
\toprule\noalign{}
\begin{minipage}[b]{\linewidth}\raggedright
Parameter
\end{minipage} & \begin{minipage}[b]{\linewidth}\raggedright
Default
\end{minipage} & \begin{minipage}[b]{\linewidth}\raggedright
Description
\end{minipage} \\
\midrule\noalign{}
\endhead
\bottomrule\noalign{}
\endlastfoot
\texttt{-\/-panel} & \texttt{Universal/Myeloid\_Panel\_Targets.chr.txt}
& Panel design file \\
\texttt{-\/-ref} & \texttt{Universal/hg19/hg19.fa} & Reference genome \\
\texttt{-\/-annotated\_panel} &
\texttt{Universal/annotated\_Myeloid\_Panel\_Targets.txt} &
Pre-annotated panel \\
\texttt{-\/-config} & \texttt{\{directory\}/config.txt} & Sample
configuration file \\
\texttt{-\/-threads} & 6 & Threads per sample \\
\texttt{-\/-parallel\_jobs} & 4 & Samples to run in parallel \\
\texttt{-\/-use\_parallel} & false & Use parallel versions of scripts
(flag) \\
\texttt{-\/-force} & false & Force re-run all stages even if outputs
exist \\
\texttt{-\/-start} & 1 & Start at this stage (1-6) \\
\texttt{-\/-stop} & 6 & Stop after this stage (1-6) \\
\end{longtable}
}

\subsection{Directory Structure}\label{directory-structure}

\begin{verbatim}
{--directory}/
├── bwa_out/                    # Stage 1 output
│   ├── *.bam
│   ├── *.bam.bai
│   ├── *.flagstat.txt
│   └── _logs/
├── filtered_bam/               # Stage 2 output
│   ├── *.filtered.bam
│   ├── *.filtered.bam.bai
│   └── _logs/
├── Read_Processing/            # Stage 3 output
│   ├── {sample}/
│   │   ├── {sample}_clean.bam
│   │   ├── {sample}_filtered_read_counts.txt
│   │   ├── {sample}_raw_coverage_per_smMIP.txt
│   │   └── {sample}_UMI_usage_per_smMIP.txt
│   ├── pileup/                 # Stage 4 output
│   │   ├── {sample}_raw_pileup.txt
│   │   └── {sample}_sscs_pileup.txt
│   └── _logs/
├── results/                    # Stage 6 output
│   └── called_mutations.txt
└── config.txt                  # Sample configuration
\end{verbatim}

\subsection{Pipeline Flow Diagram}\label{pipeline-flow-diagram}

\begin{verbatim}
FASTQ files (*.R1*.fastq.gz, *.R2*.fastq.gz)
    │
    ▼
┌─────────────────────────────────────────────────────────────┐
│ Stage 1: BWA Alignment                                      │
│ Script: BWA.sh / BWA_parallel.sh                            │
│ Input:  *R1*.fastq.gz, *R2*.fastq.gz                        │
│ Output: *.bam, *.bam.bai, *.flagstat.txt                    │
└─────────────────────────────────────────────────────────────┘
    │
    ▼
┌─────────────────────────────────────────────────────────────┐
│ Stage 2: BAM Filtering                                      │
│ Script: filter_bam.sh / filter_bam_parallel.sh              │
│ Input:  *.bam from Stage 1                                  │
│ Output: *.filtered.bam                                      │
│ Filter: MAPQ >= 50, properly paired (-f 2)                  │
└─────────────────────────────────────────────────────────────┘
    │
    ▼
┌─────────────────────────────────────────────────────────────┐
│ Stage 3: Read Processing / UMI Extraction                   │
│ Script: Read_Processing.sh / Read_Processing_parallel.sh    │
│ R Script: map_smMIPs_extract_UMIs.R                         │
│ Input:  *.filtered.bam + Panel_Targets.txt                  │
│ Output: Per sample directory containing:                    │
│         - {sample}_clean.bam                                │
│         - {sample}_filtered_read_counts.txt                 │
│         - {sample}_raw_coverage_per_smMIP.txt               │
│         - {sample}_UMI_usage_per_smMIP.txt                  │
└─────────────────────────────────────────────────────────────┘
    │
    ▼
┌─────────────────────────────────────────────────────────────┐
│ Stage 4: Pileup Generation (Level Base Calls)               │
│ Script: Level_Bas_Calls.sh / Level_Bas_Calls_parallel.sh    │
│ R Script: smMIP_level_raw_and_consensus_pileups.R           │
│ Input:  {sample}_clean.bam + Panel_Targets.txt              │
│ Output: - {sample}_raw_pileup.txt                           │
│         - {sample}_sscs_pileup.txt (consensus)              │
└─────────────────────────────────────────────────────────────┘
    │
    ▼
┌─────────────────────────────────────────────────────────────┐
│ Stage 5: Panel Annotation (ONE-TIME per panel)              │
│ R Script: Annotate_SNVs.R                                   │
│ Input:  Panel_Targets.txt                                   │
│ Output: annotated_Panel_Targets.txt                         │
│ Annotations: gene, protein, COSMIC, MAF, CADD scores        │
└─────────────────────────────────────────────────────────────┘
    │
    ▼
┌─────────────────────────────────────────────────────────────┐
│ Stage 6: Mutation Calling                                   │
│ R Script: calling_mutations.R                               │
│ Input:  - Pileup folder (*_raw_pileup.txt, *_sscs_pileup)   │
│         - Configuration file (sample info)                  │
│         - Annotated panel file                              │
│ Output: called_mutations.txt                                │
└─────────────────────────────────────────────────────────────┘
    │
    ▼
  called_mutations.txt
\end{verbatim}

\subsection{Configuration File Format}\label{configuration-file-format}

The configuration file (\texttt{config.txt}) describes sample metadata
for mutation calling.

\subsubsection{Format}\label{format}

Tab-delimited file with columns:

{\def\LTcaptype{none} % do not increment counter
\begin{longtable}[]{@{}ll@{}}
\toprule\noalign{}
Column & Description \\
\midrule\noalign{}
\endhead
\bottomrule\noalign{}
\endlastfoot
id & Sample ID (must match pileup filenames) \\
type & Sample type: \texttt{case} or \texttt{control} \\
replicate & Technical replicate pairing (or \texttt{NA}) \\
\end{longtable}
}

\subsubsection{Example}\label{example}

\begin{verbatim}
id  type    replicate
Sample1 case    NA
Sample2 case    NA
Sample3 control NA
Sample4 case    rep1
Sample5 case    rep1
\end{verbatim}

\subsubsection{Auto-Generation}\label{auto-generation}

If \texttt{-\/-config} is not provided, the pipeline will: 1. Scan the
pileup directory for \texttt{*\_raw\_pileup.txt} files 2. Extract sample
IDs from filenames 3. Generate a template \texttt{config.txt} with all
samples as \texttt{case} 4. \textbf{Stop and prompt user to edit the
file} 5. Re-run with \texttt{-\/-start\ 6} to continue

\subsection{Mutation Calling
Parameters}\label{mutation-calling-parameters}

Key parameters for \texttt{calling\_mutations.R}:

{\def\LTcaptype{none} % do not increment counter
\begin{longtable}[]{@{}lll@{}}
\toprule\noalign{}
Parameter & Default & Description \\
\midrule\noalign{}
\endhead
\bottomrule\noalign{}
\endlastfoot
\texttt{-s} & Required & Path to pileup folder \\
\texttt{-f} & Required & Configuration file (sample layout) \\
\texttt{-a} & Required & Annotated panel file \\
\texttt{-b} & ``sum'' & Binomial error rate method (``sum'' or
``max'') \\
\texttt{-g} & 10 & Min coverage for overlap (R1/R2) \\
\texttt{-m} & 0.001 & MAF cutoff to exclude known SNPs \\
\texttt{-v} & 0.05 & VAF cutoff for error modeling \\
\texttt{-p} & 0.05 & P-value cutoff (Bonferroni corrected) \\
\texttt{-d} & 2 & Fold-difference for cross-sample VAF \\
\texttt{-t} & 1 & Number of threads \\
\end{longtable}
}

\subsection{Output:
called\_mutations.txt}\label{output-called_mutations.txt}

Key columns in the final mutation calls:

{\def\LTcaptype{none} % do not increment counter
\begin{longtable}[]{@{}ll@{}}
\toprule\noalign{}
Column & Description \\
\midrule\noalign{}
\endhead
\bottomrule\noalign{}
\endlastfoot
sample\_ID & Sample identifier \\
smMIP & smMIP probe name \\
chr & Chromosome \\
pos & Position (hg19) \\
ref & Reference allele \\
alt & Alternative allele \\
gene & Gene symbol \\
protein & Protein change (e.g., p.V600E) \\
cosmic & COSMIC ID if known \\
maf & Minor allele frequency (population) \\
variant\_type & SNV, insertion, deletion \\
cadd\_scaled & CADD pathogenicity score \\
P-value & Raw p-value \\
P-value.Bonferroni & Bonferroni-corrected p-value \\
non.ref.counts & Non-reference allele counts \\
total.depth & Total read depth \\
allele.frequency & Variant allele frequency \\
SSCS.* & Consensus-based metrics \\
flags & Quality flags \\
\end{longtable}
}

\subsection{Usage Examples}\label{usage-examples}

\subsubsection{Full Pipeline Run}\label{full-pipeline-run}

\begin{Shaded}
\begin{Highlighting}[]
\FunctionTok{bash}\NormalTok{ Master\_pipeline.sh }\DataTypeTok{\textbackslash{}}
  \AttributeTok{{-}{-}directory}\NormalTok{ /path/to/experiment }\DataTypeTok{\textbackslash{}}
  \AttributeTok{{-}{-}fastq\_dir}\NormalTok{ /path/to/fastqs}
\end{Highlighting}
\end{Shaded}

\subsubsection{Parallel Processing (Recommended for Large
Datasets)}\label{parallel-processing-recommended-for-large-datasets}

\begin{Shaded}
\begin{Highlighting}[]
\FunctionTok{bash}\NormalTok{ Master\_pipeline.sh }\DataTypeTok{\textbackslash{}}
  \AttributeTok{{-}{-}directory}\NormalTok{ /path/to/experiment }\DataTypeTok{\textbackslash{}}
  \AttributeTok{{-}{-}fastq\_dir}\NormalTok{ /path/to/fastqs }\DataTypeTok{\textbackslash{}}
  \AttributeTok{{-}{-}use\_parallel} \DataTypeTok{\textbackslash{}}
  \AttributeTok{{-}{-}parallel\_jobs}\NormalTok{ 4 }\DataTypeTok{\textbackslash{}}
  \AttributeTok{{-}{-}threads}\NormalTok{ 6}
\end{Highlighting}
\end{Shaded}

\subsubsection{Resume from Specific
Stage}\label{resume-from-specific-stage}

\begin{Shaded}
\begin{Highlighting}[]
\CommentTok{\# After editing config.txt, run only mutation calling}
\FunctionTok{bash}\NormalTok{ Master\_pipeline.sh }\DataTypeTok{\textbackslash{}}
  \AttributeTok{{-}{-}directory}\NormalTok{ /path/to/experiment }\DataTypeTok{\textbackslash{}}
  \AttributeTok{{-}{-}fastq\_dir}\NormalTok{ /path/to/fastqs }\DataTypeTok{\textbackslash{}}
  \AttributeTok{{-}{-}start}\NormalTok{ 6 }\AttributeTok{{-}{-}stop}\NormalTok{ 6}
\end{Highlighting}
\end{Shaded}

\subsubsection{Force Re-run All Stages}\label{force-re-run-all-stages}

\begin{Shaded}
\begin{Highlighting}[]
\FunctionTok{bash}\NormalTok{ Master\_pipeline.sh }\DataTypeTok{\textbackslash{}}
  \AttributeTok{{-}{-}directory}\NormalTok{ /path/to/experiment }\DataTypeTok{\textbackslash{}}
  \AttributeTok{{-}{-}fastq\_dir}\NormalTok{ /path/to/fastqs }\DataTypeTok{\textbackslash{}}
  \AttributeTok{{-}{-}force}
\end{Highlighting}
\end{Shaded}

\subsubsection{Custom Panel and
Reference}\label{custom-panel-and-reference}

\begin{Shaded}
\begin{Highlighting}[]
\FunctionTok{bash}\NormalTok{ Master\_pipeline.sh }\DataTypeTok{\textbackslash{}}
  \AttributeTok{{-}{-}directory}\NormalTok{ /path/to/experiment }\DataTypeTok{\textbackslash{}}
  \AttributeTok{{-}{-}fastq\_dir}\NormalTok{ /path/to/fastqs }\DataTypeTok{\textbackslash{}}
  \AttributeTok{{-}{-}panel}\NormalTok{ /path/to/custom\_panel.txt }\DataTypeTok{\textbackslash{}}
  \AttributeTok{{-}{-}ref}\NormalTok{ /path/to/hg38.fa }\DataTypeTok{\textbackslash{}}
  \AttributeTok{{-}{-}annotated\_panel}\NormalTok{ /path/to/annotated\_panel.txt}
\end{Highlighting}
\end{Shaded}

\subsection{Resource Requirements}\label{resource-requirements}

\subsubsection{Memory}\label{memory}

{\def\LTcaptype{none} % do not increment counter
\begin{longtable}[]{@{}lll@{}}
\toprule\noalign{}
Stage & Memory per Sample & Notes \\
\midrule\noalign{}
\endhead
\bottomrule\noalign{}
\endlastfoot
1 (BWA) & \textasciitilde8 GB & Scales with reference genome size \\
2 (Filter) & \textasciitilde4 GB & I/O bound \\
3 (ReadProc) & \textasciitilde30 GB & R memory-intensive \\
4 (Pileup) & \textasciitilde25 GB & R memory-intensive \\
5 (Annotate) & \textasciitilde8 GB & Network-bound (API calls) \\
6 (Mutations) & \textasciitilde16 GB & Depends on cohort size \\
\end{longtable}
}

\subsubsection{Recommended Settings}\label{recommended-settings}

{\def\LTcaptype{none} % do not increment counter
\begin{longtable}[]{@{}lll@{}}
\toprule\noalign{}
System RAM & --parallel\_jobs & --threads \\
\midrule\noalign{}
\endhead
\bottomrule\noalign{}
\endlastfoot
32 GB & 2 & 4 \\
64 GB & 3 & 6 \\
128 GB & 4 & 6 \\
256 GB & 8 & 6 \\
\end{longtable}
}

\subsection{Dependencies}\label{dependencies}

\subsubsection{Command-Line Tools}\label{command-line-tools}

\begin{itemize}
\tightlist
\item
  bwa (or bwa-mem2)
\item
  samtools
\item
  GNU parallel (for --use\_parallel)
\end{itemize}

\subsubsection{R Packages}\label{r-packages}

\begin{itemize}
\tightlist
\item
  optparse
\item
  data.table
\item
  parallel
\item
  cellbaseR (for annotation)
\item
  IRanges
\item
  Rsamtools
\end{itemize}

\subsection{Troubleshooting}\label{troubleshooting}

\subsubsection{Out of Memory}\label{out-of-memory}

\begin{itemize}
\tightlist
\item
  Reduce \texttt{-\/-parallel\_jobs}
\item
  Ensure R\_MAX\_VSIZE is set (automatically handled by pipeline)
\end{itemize}

\subsubsection{Missing R1/R2 Pairs}\label{missing-r1r2-pairs}

\begin{itemize}
\tightlist
\item
  Check FASTQ naming convention matches \texttt{*R1*.fastq.gz} pattern
\item
  Use \texttt{-\/-r1\_glob} and \texttt{-\/-r2\_token} to customize
  pattern
\end{itemize}

\subsubsection{Stage Skip Not Working}\label{stage-skip-not-working}

\begin{itemize}
\tightlist
\item
  Use \texttt{-\/-force} to override skip logic
\item
  Check that expected output files exist and are non-empty
\end{itemize}

\subsubsection{Config File Issues}\label{config-file-issues}

\begin{itemize}
\tightlist
\item
  Ensure tab-delimited format
\item
  Sample IDs must match pileup filenames exactly
\item
  All samples need a type assignment (case or control)
\end{itemize}

\end{document}
